% ============================================================================ 3 · Definición del desafío
\section{Definición del desafío}
\label{sec:definicion}

\subsection{Desafío específico: anticipar el momento, sin hacer trampa}

\textbf{Qué entendimos.} Ford pide un modelo que anticipe, a partir de patrones históricos de uso y
variables de motor, \emph{el momento en que un vehículo se aproxima a un evento de degradación}, a tiempo
para una acción preventiva antes de que el cliente perciba el problema.

Esa frase tiene dos exigencias que tiran para lados distintos. \emph{Anticipar} pide no mirar los últimos
kilómetros antes de la falla, que es donde está la señal más fuerte. \emph{El momento} pide saber
\emph{cuándo}, y no solo \emph{qué auto}. La mayor parte del trabajo metodológico fue convertir esas dos
exigencias en una pregunta medible que no se pudiera ganar con atajos.

\subsubsection{La pregunta: un panel de cortes con gap de blanking}
\label{sec:panel}

Armamos un panel con una fila por par (vehículo, punto de corte). Los cortes caen cada
\textbf{Δ = 500 km} de odómetro. En cada corte, el modelo ve solo los últimos \textbf{W = 1.000 km} y
responde: ¿el evento cae entre \textbf{G = 500} y \textbf{G + H = 3.500 km} después del corte?
(figura~\ref{fig:panel}).

\begin{figure}[H]
\centering
\begin{tikzpicture}[x=1.55cm,y=1cm,font=\sffamily\scriptsize]
\fill[fordblue!18] (0,0) rectangle (1,0.7);
\fill[muted!35] (1,0) rectangle (1.5,0.7);
\fill[accent2!22] (1.5,0) rectangle (4.5,0.7);
\draw[inksec] (-0.6,0) -- (5.1,0);
\foreach \x/\l in {0/{−1.000},1/{corte},1.5/{+500},4.5/{+3.500}} {\draw[inksec] (\x,0) -- (\x,-0.1) node[below] {\l};}
\node at (0.5,0.35) {\textbf{W} · lo que mira};
\node at (1.25,0.35) {\textbf{G}};
\node at (3,0.35) {\textbf{H} · ¿falla acá?};
\draw[very thick,ink] (1,-0.15) -- (1,0.95);
\node[above] at (1,0.95) {punto de corte};
\node[below,text width=3.2cm,align=center] at (1.25,-0.45) {500 km ciegos: el modelo nunca los ve};
\node[right] at (5.1,0) {km};
\end{tikzpicture}
\caption{El planteo: ventana W hacia atrás, gap G y horizonte H hacia adelante. La etiqueta vale 1 si el
evento cae en el horizonte.}
\label{fig:panel}
\fuente{\repo{configs/\allowbreak{}data/\allowbreak{}panel\_v2.yaml}; \repo{docs/\allowbreak{}memoria/\allowbreak{}decisiones.md} (18-09).}
\end{figure}

\begin{itemize}
  \item \textbf{G = 500 km} es la regla antirreactiva. Sin ese gap, un modelo que mire los últimos kilómetros
    antes del evento memoriza los mensajes de “filtro sobrecargado” y reporta un número alto haciendo
    detección reactiva.
  \item \textbf{W = 1.000 km.} Con la primera entrega, cada 500 km más de ventana costaba unos tres eventos
    con positivo, y ventanas de 1.000, 2.000 y 3.000 km daban la misma métrica dentro del ruido.
  \item \textbf{H = 3.000 km} pone los positivos entre 500 y 3.500 km antes del evento, donde el perfil
    alineado muestra que la señal existe (§\ref{sec:eda}).
  \item \textbf{El eje es el odómetro} para cortar las ventanas, pero el evento se ordena por tiempo y no por
    km: la dispersión del logaritmo es 0,38 en edad contra 0,91 en odómetro. Quien maneja más llega al
    evento con más km, no antes. Por eso la anticipación se reporta en km y también en días aproximados.
  \item \textbf{Sanos emparejados por odómetro y por mes.} Los eventos se concentran en ciertos meses. Sin
    emparejar por mes, sumar tres columnas de calendario subía el ROC de 0,67 a 0,76; emparejando, de 0,57
    a 0,60 (primera entrega). El calendario no es el auto.
\end{itemize}

Con la entrega v2, el panel de desarrollo tiene 11.703 filas de 426 autos con cortes (135 fallados, 291
sanos) y 682 filas positivas: una tasa base de 0,058 por fila. Los autos sanos aportan en promedio 26,9
cortes y los fallados 28,8: las bolsas quedaron casi simétricas, algo que con la primera entrega no pasaba
(9,0 contra 18,3) y que obligó a corregir las métricas (abajo).

\subsubsection{Las reglas que no se negocian}

\begin{enumerate}
  \item \textbf{Gap de blanking}: el modelo nunca ve los 500 km previos al evento.
  \item \textbf{Split agrupado por vehículo}: un auto nunca está a la vez en entrenamiento y en validación.
    La lógica vive en un solo módulo (\repo{src/\allowbreak{}eval/\allowbreak{}splits.py}). Los folds se estratifican por la etiqueta a
    nivel vehículo, cruzada con mercado y motor, y la CV se repite tres veces.
  \item \textbf{Holdout congelado}: 111 autos sorteados una vez, antes de mirar ninguna feature contra la
    etiqueta. Todo lo que se mira es desarrollo; las filas de test no se tocan hasta tener el modelo elegido.
  \item \textbf{Features solo hacia atrás}: nada posterior al corte. Escalado e imputación se ajustan con el
    train de cada fold, dentro del pipeline.
  \item \textbf{Auditorías obligatorias} antes de aceptar un número (§\ref{sec:metricas}, Anexo~\ref{anx:auditorias}).
  \item \textbf{Nada se hardcodea}: paths, semillas e hiperparámetros salen de un YAML versionado.
\end{enumerate}

\subsubsection{Cómo medimos, y por qué cambiamos de métrica}
\label{sec:metricas}

Empezamos midiendo como manda el libro: PR-AUC por fila. Los primeros modelos daban números que parecían
buenos, hasta que nos preguntamos qué pasaría con un modelo que solo reconoce \emph{qué auto} falla, sin
saber nada del \emph{cuándo}.

\begin{caja}[title={El techo de cohorte}]
\small
Todas las filas positivas del panel están dentro de los cortes de autos que fallan. Un modelo que
puntúa cada fila con “¿este auto es de los que fallan?”, sin ninguna información del \emph{cuándo}, saca
\textbf{PR-AUC 0,177 y lift 3,03×} en el dev v2 (0,263 y 2,10× con la primera entrega). La meta inicial del
plan (lift de 1,6–2×) se alcanzaba sin anticipar nada. Ningún modelo superó ese techo: \textbf{un PR-AUC por
fila de este panel mide sobre todo \emph{qué auto} y muy poco \emph{cuándo}.}
\fuente{\repo{docs/\allowbreak{}memoria/\allowbreak{}decisiones.md}, cierre de F3 (20-09); \repo{src/\allowbreak{}eval/\allowbreak{}metrics.py::cohort\_ceiling}.}
\end{caja}

Tres mediciones independientes llegaron al mismo lugar con la primera entrega: colapsar el puntaje de cada
auto a su promedio \emph{mejoraba} la decisión por vehículo; la auditoría \aprime{} del modelo de control
daba negativa en las tres repeticiones; y un nulo que permuta la etiqueta dentro de cada auto quedaba por
encima del modelo. Desde entonces la pregunta que decide es una sola:

\begin{quote}
\emph{De los autos que van a fallar, ¿a cuántos les avisamos a tiempo si solo permitimos una cantidad fija
de falsas alarmas entre los autos sanos?}
\end{quote}

\begin{itemize}
  \item \textbf{Detección por auto} con umbral exacto: para cada presupuesto b (2 a 30\,\% de los sanos con
    alguna falsa alarma), el umbral es el menor que deja a lo sumo b de los sanos alertados. Una alerta son
    \textbf{dos cortes seguidos} sobre el umbral. El criterio primario es la detección media al 5, 10, 15
    y 20\,\%.
  \item \textbf{Contra el azar, con el mismo historial}: el puntaje permutado entre filas, 200 veces. Con la
    primera entrega, las bolsas asimétricas “detectaban” 10\,\% al 5\,\% sin ninguna señal; toda métrica que
    dependa del largo del historial se compara contra un nulo que lo conserve.
  \item \textbf{Contra la celda mercado × motor}: la tasa de fallas de la celda del auto, calculada sin él.
    Tres celdas (Chile × ENG\_2, Colombia × ENG\_2 y Brasil × ENG\_3) concentran 117 de los 135 fallados de
    dev. Esa regla, sin mirar un viaje, detecta 18 · 33 · 62\,\% al 5 · 10 · 20\,\% \tagdev{}.
  \item \textbf{Fuera de muestra}: el umbral se fija con los sanos de los otros folds y se mide la tasa de
    falsas alarmas que realmente produce.
  \item \textbf{Bootstrap pareado por vehículo} para comparar dos modelos, con 1.000 a 2.000 réplicas.
  \item \textbf{AUC dentro de mercado × motor}: solo compara pares de autos de la misma celda. La diferencia
    con el AUC agrupado es composición de la muestra.
\end{itemize}

Para elegir candidatos, además, cada comparación nueva gastaba un lugar de un \emph{presupuesto de
comparaciones}: con unos 12 eventos por fold, sumar candidatos sobre la marcha garantiza que “el mejor” sea
ruido. Las variantes importantes se preregistraron por escrito antes de medirlas (\repo{docs/\allowbreak{}memoria/\allowbreak{}*preregistro*}).

\subsection{Impacto operativo}

\textbf{Qué entendimos.} No anticipar estos eventos afecta la satisfacción del cliente, sube los costos de
garantía y reparación, reduce la vida útil percibida y desordena la planificación de mantenimiento de la red
de concesionarios.

\textbf{Qué hicimos con cada punto.}
\begin{itemize}
  \item \textbf{Satisfacción}: un aviso a tiempo, en lenguaje humano y con algo concreto para hacer. Si el
    auto no iba a fallar, el conductor igual recibió un consejo que le sirve.
  \item \textbf{Garantía y reparación}: la alerta llega con una mediana de unos tres meses de margen
    \tagtest{}, tiempo para cambiar un hábito o programar un turno en vez de una urgencia. Cuántas fallas se
    evitan no se puede medir con estos datos: es la métrica central del piloto (§\ref{sec:viabilidad}).
  \item \textbf{Planificación de concesionarios}: la bandeja semanal prioriza los casos, y la perilla deja que
    Ford elija cuántas alertas acepta según su capacidad y el costo de cada una.
\end{itemize}
